\section{Objet et portée de l'annexe}
\label{annexe-I-objet}

Cette annexe établit la suffisance de l'état lent aligné pour la réponse
à une sonde déclarée. Elle prolonge l'Appendice~H : celui-ci traitait la
section éteinte $p=\chi=\theta=0$ et montrait que le contenu survivant
$\mathbf s$ suffit alors à prédire l'issue. Ici, les variables lentes
résiduelles sont conservées. Le résumé devient
\[
z=(p,\chi,\theta,s),
\]
sur une famille où les variables vectorielles internes demeurent
alignées sur une direction déclarée $\mathbf u_1$.

Le résultat est une proposition de fermeture, non un principe universel.
Il dépend de trois faits à vérifier : l'invariance du sous-espace aligné,
l'inertie de toute variable omise et le déterminisme de l'exécution. Il
ne prétend ni que quatre coordonnées suffisent hors alignement, ni que
la même réduction subsiste sous une loi qui lirait une variable
actuellement inerte.

\begin{encadre}[title={Résultat directeur}]
Sur la famille alignée et sous la loi $\mathcal G_{\mathrm{lin}}$, la
réponse à une sonde admissible $\mathcal P$ est exactement paramétrée
par l'état lent
\[
z=(p,\chi,\theta,s).
\]
Il existe une carte
\[
\widehat\Omega_{\mathcal P,\rho}:Z\longrightarrow\mathcal Y_\rho
\]
telle que
\[
\mathcal R_z^\rho(\mathcal P)
=
\widehat\Omega_{\mathcal P,\rho}(p,\chi,\theta,s).
\]
La réduction de l'Appendice~H est la restriction de cette carte à la
section $p=\chi=\theta=0$.
\end{encadre}

\section{Cadre déclaré}
\label{annexe-I-cadre}

\subsection{État complet et regraine canonique}
\label{annexe-I-etat-complet}

On note l'état complet du modèle canonique
\[
X=(y,p,\boldsymbol\chi,\boldsymbol\theta,\mathbf s,ssc),
\]
où $y$ rassemble les variables mécaniques, $p$ la pression scalaire,
$\boldsymbol\chi$ la variable d'interface, $\boldsymbol\theta$ la phase
intermédiaire, $\mathbf s$ le contenu mnésique et $ssc$ une variable
interne additionnelle du modèle de référence.

La regraine canonique $\mathcal G$ appartient au protocole, conformément
à A-H1. Elle impose le même état mécanique $y_0$ aux systèmes comparés
et préserve les variables internes :
\[
\Phi_{\mathcal G}
(y,p,\boldsymbol\chi,\boldsymbol\theta,\mathbf s,ssc)
=
(y_0,p,\boldsymbol\chi,\boldsymbol\theta,\mathbf s,ssc).
\]
La sonde commence donc depuis une section mécanique commune, sans effacer
les traces lentes dont on cherche l'effet.

\subsection{Sous-espace aligné}
\label{annexe-I-sous-espace-aligne}

Fixons une direction unitaire $\mathbf u_1$ appartenant au sous-espace des
points fixes d'une équivariance déclarée du flot et des événements. Dans
le modèle canonique, on choisit la diagonale
$\mathbf u_1=(1,1)/\sqrt2$, fixe par échange des composantes. Les axes
et l'anti-diagonale peuvent également être traités lorsqu'ils sont
préservés par les entrées et les règles considérées. Une direction
générique n'est pas supposée admissible. Le sous-espace aligné est
l'ensemble
\[
\mathcal M_{\mathrm{al}}
=
\left\{
(y,p,\boldsymbol\chi,\boldsymbol\theta,\mathbf s,ssc)
\ \middle|\
\boldsymbol\chi=\chi\mathbf u_1,
\ \boldsymbol\theta=\theta\mathbf u_1,
\ \mathbf s=s\mathbf u_1
\right\}.
\]
Après regraine, la partie pertinente de cet état est décrite par
\[
z=(p,\chi,\theta,s)\in Z\subseteq\mathbb R^4.
\]
On définit l'injection canonique
\begin{equation}
\label{eq:annexe-I-injection}
\iota(z;ssc)
=
(y_0,p,\chi\mathbf u_1,\theta\mathbf u_1,s\mathbf u_1,ssc).
\end{equation}

Le paramètre $ssc$ est laissé visible dans
\eqref{eq:annexe-I-injection} afin de ne pas confondre égalité littérale
des états et suffisance comportementale du résumé $z$.

\subsection{Lecture et segmentation}
\label{annexe-I-lecture}

Le couple d'observation $(\mathfrak B,\rho)$ et la sonde $\mathcal P$
sont déclarés. Pour une lecture terminale, aucune condition
supplémentaire n'est nécessaire. Pour une lecture de chemin, on suppose
la compatibilité à la segmentation et au suffixe définie dans
l'Appendice~G. La réponse est
\[
\mathcal R_X^\rho(\mathcal P)
=
\omega_\rho\!\left(
\operatorname{Exec}_{\mathcal P}(X)
\right).
\]

\section{Hypothèses de fermeture}
\label{annexe-I-hypotheses}

\paragraph{(A1) --- Déterminisme.}
Pour tout $X$ de la classe considérée et toute sonde admissible
$\mathcal P$, l'exécution $\operatorname{Exec}_{\mathcal P}(X)$ est
unique.

\paragraph{(A2) --- Invariance du sous-espace aligné.}
Le flot continu, les événements et les épisodes de la sonde préservent
$\mathcal M_{\mathrm{al}}$. Autrement dit, un état aligné soumis à des
entrées alignées ne produit aucune composante transverse.

\paragraph{(A3) --- Suffisance comportementale de $z$.}
Pour deux états alignés ayant le même quadruplet $z$, toute composante
interne omise du résumé est comportementalement inerte sous la sonde et
la loi déclarées. Formellement, pour tous $ssc_1,ssc_2$ admissibles,
\begin{equation}
\label{eq:annexe-I-inertie}
\mathcal R_{\iota(z;ssc_1)}^\rho(\mathcal P)
=
\mathcal R_{\iota(z;ssc_2)}^\rho(\mathcal P).
\end{equation}

\paragraph{(A4) --- Regraine commune.}
Les systèmes comparés sont ramenés au même état mécanique autonome
$y_0$ avant la sonde, sans altération de $z$.

\begin{encadre}[title={Statut de l'hypothèse markovienne}]
(A3) est l'hypothèse décisive. Elle ne découle pas de l'alignement. Dans
le modèle canonique sous $\mathcal G_{\mathrm{lin}}$, $ssc$ n'apparaît
dans aucune équation active : elle est donc inerte et
\eqref{eq:annexe-I-inertie} est satisfaite. Sous une loi lisant $ssc$,
le résumé devrait devenir $(p,\chi,\theta,s,ssc)$, ou la proposition
cesserait d'être applicable.
\end{encadre}

\section{Lemme d'invariance alignée}
\label{annexe-I-lemme-invariance}

\resultstep{lem:I-invariance-alignee}
\begin{encadre}[title={Lemme \theproposition{} --- Invariance de la famille alignée}]
Supposons qu'une symétrie déclarée du modèle fixe $\mathbf u_1$, que le
flot et les règles événementielles soient équivariants sous cette
symétrie, et que la graine, les entrées et les contenus appartiennent à
son sous-espace de points fixes. Alors
$\mathcal M_{\mathrm{al}}$ est invariant par toute exécution admissible
de la sonde.
\end{encadre}

\paragraph{Démonstration.}
Écrivons toute variable vectorielle $q$ comme
\[
q=q_\parallel\mathbf u_1+q_\perp,
\qquad q_\perp\cdot\mathbf u_1=0.
\]
Par équivariance, le champ de vecteurs restreint à $q_\perp=0$ ne crée
aucune composante transverse : la dérivée de $q_\perp$ y est nulle. Par
unicité du flot, une trajectoire initialement dans cette sous-variété y
reste jusqu'au premier événement. Les règles événementielles appliquent
les mêmes opérations aux composantes symétriques et reçoivent des
entrées alignées ; elles préservent donc elles aussi $q_\perp=0$. Une
induction sur la suite finie des épisodes et événements établit
l'invariance de l'exécution entière. \hfill$\square$

\paragraph{Portée.}
Le lemme ne dit rien des perturbations transverses ni de leur stabilité.
Il établit une fermeture exacte de la famille déclarée. La stabilité au
voisinage de $\mathcal M_{\mathrm{al}}$ est une question distincte.

\section{Suffisance de l'état lent aligné}
\label{annexe-I-suffisance-etat-lent}

\resultstep{prop:I-etat-lent-aligne}
\begin{encadre}[title={Proposition \theproposition{} --- Paramétrage exact par l'état lent aligné}]
Sous (A1)--(A4), pour toute sonde admissible $\mathcal P$ et toute
résolution $\rho$, il existe une application unique sur le domaine
atteignable
\[
\widehat\Omega_{\mathcal P,\rho}:Z\longrightarrow\mathcal Y_\rho
\]
telle que, pour tout état aligné regrainé de résumé
$z=(p,\chi,\theta,s)$,
\begin{equation}
\label{eq:annexe-I-proposition}
\boxed{
\mathcal R_z^\rho(\mathcal P)
=
\widehat\Omega_{\mathcal P,\rho}(p,\chi,\theta,s).
}
\end{equation}
Deux histoires qui déposent le même $z$ avant la sonde sont donc
indiscernables par cette sonde à la résolution $\rho$.
\end{encadre}

\paragraph{Démonstration.}
Par (A4), tout état comparé possède la même composante mécanique $y_0$.
Par le Lemme~\ref{lem:I-invariance-alignee} et (A2), l'exécution reste dans
$\mathcal M_{\mathrm{al}}$ ; ses composantes vectorielles actives sont
donc entièrement décrites par les scalaires $(\chi,\theta,s)$, auxquels
s'ajoute $p$. Pour $z\in Z$, choisissons une valeur admissible quelconque
$ssc_0$ et posons
\begin{equation}
\label{eq:annexe-I-definition-Omega}
\widehat\Omega_{\mathcal P,\rho}(z)
:=
\omega_\rho\!\left(
\operatorname{Exec}_{\mathcal P}
\bigl(\iota(z;ssc_0)\bigr)
\right).
\end{equation}
Par (A3), cette définition ne dépend pas du choix de $ssc_0$. Par (A1),
l'exécution et sa lecture sont univoques. Ainsi
$\widehat\Omega_{\mathcal P,\rho}$ est bien définie et toute histoire
qui dépose le même $z$ reçoit la même réponse. L'unicité sur le domaine
atteignable est immédiate, puisque la valeur de la carte y est imposée
par la réponse du système. \hfill$\square$

\section{Relation avec la réduction au contenu survivant}
\label{annexe-I-relation-H}

Soit la section éteinte
\[
Z_{\mathrm{ext}}
=
\{(0,0,0,s):s\in I\}.
\]
La carte nue de l'Appendice~H est la restriction de la carte lente :
\begin{equation}
\label{eq:annexe-I-restriction-H}
\boxed{
\Omega_{\mathcal P,\rho}(s)
=
\widehat\Omega_{\mathcal P,\rho}(0,0,0,s).
}
\end{equation}
Le Théorème~\ref{thm:H-contenu-survivant} n'est donc pas un mécanisme séparé. Il décrit la trace
de la carte complète sur la section où toutes les autres coordonnées
lentes sont éteintes.

\resultstep{cor:I-emboitement-resumes}
\begin{encadre}[title={Corollaire \theproposition{} --- Emboîtement des résumés suffisants}]
Si un protocole d'oubli annule successivement certaines coordonnées de
$z$, la réponse se factorise par le résumé restant. En particulier :
\[
(p,\chi,\theta,s)
\longrightarrow
(\theta,s)
\longrightarrow
s
\]
chaque fois que les coordonnées supprimées sont exactement éteintes ou
comportementalement inertes sous la sonde déclarée. Sous l'extension
plastique, la même mécanique conduit au cinquième étage
$z_5=(p,\chi,\theta,s,\gamma)$ étudié dans l'Appendice~\ref{annexe-equivalence-presentation-plasticite}.
\end{encadre}

Cette chaîne n'est pas une hiérarchie ontologique des variables. Elle
exprime une hiérarchie expérimentale de résumés suffisants, relative au
protocole et à la résolution.

\section{Ensemble d'appropriation et frontière}
\label{annexe-I-frontiere}

Pour une lecture binaire, définissons
\[
A_{\mathcal P,\rho}
=
\left\{
z\in Z
\ \middle|\
\widehat\Omega_{\mathcal P,\rho}(z)=1
\right\}.
\]
Sa frontière relative dans $Z$ est
\begin{equation}
\label{eq:annexe-I-Sigma}
\Sigma_{\mathcal P,\rho}
=
\partial_Z A_{\mathcal P,\rho}.
\end{equation}
Elle rassemble les états lents pour lesquels une variation arbitrairement
petite, à la résolution déclarée, peut modifier l'issue.

\resultstep{cor:I-localisation-bascules}
\begin{encadre}[title={Corollaire \theproposition{} --- Localisation des bascules d'une famille historique}]
Soit $H:I\to Z$ la carte qui associe à un paramètre initial $a$ l'état
lent déposé par une histoire complète avant la sonde. Sous les
hypothèses de la Proposition~\ref{prop:I-etat-lent-aligne}, les points de bascule de la famille
sont contenus dans
\[
H^{-1}(\Sigma_{\mathcal P,\rho}).
\]
Si $H$ est continue et si la lecture est localement constante sur chaque
composante de $Z\setminus\Sigma_{\mathcal P,\rho}$, alors tout changement effectif d'issue se produit en un tel antécédent.
\end{encadre}

\paragraph{Démonstration.}
La réponse de la famille est
$\widehat\Omega_{\mathcal P,\rho}\circ H$. Tant que $H(a)$ demeure dans
une composante où la carte est localement constante, l'issue ne change
pas. Une bascule exige donc que la trajectoire historique rencontre la
frontière. Sous les hypothèses additionnelles, tout passage effectif
d'une composante de réponse à une autre traverse cette frontière.
\hfill$\square$

La section étudiée dans l'Appendice~H vérifie
\[
\partial A^{\mathrm{ext}}_{\mathcal P,\rho}
=
\Sigma_{\mathcal P,\rho}
\cap Z_{\mathrm{ext}}
\]
lorsque l'intersection est transverse et qu'aucune portion de frontière
ne contient la section entière. En général, la frontière de la
restriction est à comprendre relativement à $Z_{\mathrm{ext}}$ ; elle
peut différer de la simple intersection en cas de tangence ou de
coïncidence locale.

\section{Activité des coordonnées lentes}
\label{annexe-I-activite}

La suffisance d'un résumé n'implique pas que chacune de ses coordonnées
soit active pour toute sonde. Il faut distinguer la structure
géométrique portée par une coordonnée de son activité décisionnelle à la
résolution déclarée.

Pour une coordonnée $z_i$, notons $\operatorname{abl}_i(z)$ l'état obtenu
par l'intervention déclarée d'ablation, généralement $z_i:=0$, les autres
coordonnées étant conservées. Cette intervention est un protocole de
test privilégié ; elle n'est pas la définition générale de l'activité,
qui peut être établie par toute paire ne différant que selon $z_i$.

\paragraph{Activité ponctuelle.}
La coordonnée $z_i$ est active au point $z$ lorsque
\[
\widehat\Omega_{\mathcal P,\rho}(z)
\ne
\widehat\Omega_{\mathcal P,\rho}(\operatorname{abl}_i(z)).
\]
Cette notion dépend du point et de l'intervention d'ablation choisie.

\paragraph{Activité locale.}
La coordonnée $z_i$ est active localement au voisinage de $z$ s'il existe
un voisinage $U\ni z$ sur lequel la carte de réponse ne se factorise pas
par la projection oubliant $z_i$. Équivalemment, il existe dans tout
voisinage suffisamment petit une paire de points ne différant que selon
$z_i$ et donnant des réponses distinctes.

\paragraph{Activité globale sur une classe.}
Pour une classe déclarée $D\subseteq Z$, la coordonnée $z_i$ est active
globalement sur $D$ s'il existe $z,z'\in D$ ne différant que selon
$z_i$ tels que
\[
\widehat\Omega_{\mathcal P,\rho}(z)
\ne
\widehat\Omega_{\mathcal P,\rho}(z').
\]
Elle est globalement inactive sur $D$ si aucune telle paire n'existe.

Ces trois niveaux ne se confondent pas. Une coordonnée peut être inactive
au point testé, active dans un voisinage plus large, ou n'être révélée
que par une autre classe de sondes. Pour les tableaux de synthèse du
Tome III, le codage est : $0$ pour aucune activité détectée, $L$ pour
activité locale, et $P+L$ lorsqu'une activité ponctuelle par ablation est
également observée.

\resultstep{cor:I-factorisation-ablation}
\begin{encadre}[title={Corollaire \theproposition{} --- Factorisation et test par ablation}]
Si $z_i$ est globalement inactive sur le domaine atteignable $D$, alors
la réponse se factorise par la projection qui oublie $z_i$. Réciproquement,
toute paire séparatrice ne différant que selon $z_i$ interdit cette
factorisation. L'ablation fournit une telle paire lorsqu'elle conserve
les deux états dans $D$ ; elle est donc un cas d'usage suffisant, mais
non la définition exhaustive du critère.
\end{encadre}

Ce critère fournit deux voies distinctes. L'insuffisance d'un résumé se
démontre par exhibition d'une paire séparatrice. Sa suffisance exige une
preuve de fermeture, ou une vérification exhaustive sur un domaine fini
déclaré ; l'absence de séparation dans un échantillon ne suffit pas.

\section{Vérification dans le modèle canonique}
\label{annexe-I-verification-canonique}

La Proposition~\ref{prop:I-etat-lent-aligne} s'applique au modèle de référence sous les
conditions suivantes.

\paragraph{Symétrie et alignement.}
Les équations sont symétriques sous l'échange des composantes. La graine,
les drives et les contenus utilisés dans la famille de référence sont
alignés. Le sous-espace
\[
x_1=x_2,
\quad v_1=v_2,
\quad \theta_1=\theta_2,
\quad \chi_1=\chi_2,
\quad s_1=s_2
\]
est donc invariant par le flot et les événements.

\paragraph{Inertie de $ssc$.}
Sous $\mathcal G_{\mathrm{lin}}$, $ssc$ n'intervient dans aucune équation
active. Deux états ayant le même $z$ et des valeurs différentes de
$ssc$ engendrent la même exécution observable. Cette vérification est
spécifique à la loi déclarée.

\paragraph{Déterminisme.}
Le modèle hybride satisfait O1 sur les protocoles admissibles considérés.
Les règles de saut et le flot déterminent donc une exécution unique.

\paragraph{Fermeture numérique.}
Les états de regraine produits par les histoires-témoins ont été
réinjectés à partir de leur quadruplet synthétique. Les résidus
transverses observés sont nuls sur la famille déclarée et les réponses
coïncident cellule à cellule. Cette vérification ne remplace pas la
preuve d'invariance ; elle contrôle que l'implémentation numérique suit
bien le sous-modèle démontré.

\section{Ce que la proposition ne suppose pas}
\label{annexe-I-non-hypotheses}

La Proposition~\ref{prop:I-etat-lent-aligne} ne requiert pas :
\begin{itemize}
\item que $\widehat\Omega_{\mathcal P,\rho}$ soit continue ;
\item que $A_{\mathcal P,\rho}$ soit connexe ;
\item que $\Sigma_{\mathcal P,\rho}$ soit une variété lisse ;
\item que la réponse soit monotone en $s$, $\theta$, $\chi$ ou $p$ ;
\item que toutes les coordonnées de $z$ soient actives ;
\item que le résumé demeure suffisant pour une autre sonde, une autre
  résolution ou une autre loi de mémoire.
\end{itemize}

Elle établit uniquement une factorisation exacte de la réponse dans le
cadre déclaré.

\section{Limites de portée}
\label{annexe-I-limites}

\paragraph{Hors alignement.}
Si $\boldsymbol\chi$, $\boldsymbol\theta$ et $\mathbf s$ ne sont pas
colinéaires, les quatre scalaires $(p,\chi,\theta,s)$ ne déterminent plus
les directions relatives. Le résumé doit être remplacé par l'état lent
vectoriel complet, ou par tout autre quotient dont la suffisance est
démontrée.

\paragraph{Loi lisant $ssc$.}
Dès que $ssc$ intervient dans le flot, les événements ou la lecture,
(A3) tombe. Deux histoires de même $z$ peuvent alors être séparées. Il
faut ajouter $ssc$ au résumé.

\paragraph{Regraine non canonique.}
Si la regraine dépend de l'histoire ou ne ramène pas les systèmes au
même état mécanique autonome, $z$ ne suffit plus à déterminer l'état de
départ de la sonde.

\paragraph{Entrées transverses.}
Une sonde non alignée peut quitter $\mathcal M_{\mathrm{al}}$ même si
l'état initial y appartient. L'invariance doit alors être redémontrée
dans un espace plus grand.

\paragraph{Lectures historiques intégrales.}
Pour une lecture portant sur le chemin entier, la simple égalité des
états terminaux ne suffit pas à comparer des histoires complètes. La
factorisation de cette annexe concerne la réponse de la sonde suffixe à
partir de l'état regrainé ; son transport à une histoire concaténée
requiert la compatibilité au suffixe de l'Appendice~G.

\paragraph{Extinction approchée.}
Lorsque certaines coordonnées sont petites mais non nulles, on ne peut
les supprimer sans hypothèse quantitative de stabilité. Pour une observable pré-binaire, une borne de la
forme
\[
d_\rho\!\left(
\widehat\Omega(p,\chi,\theta,s),
\Omega(s)
\right)
\le
L_p|p|+L_\chi|\chi|+L_\theta|\theta|
\]
exige une régularité locale qui n'est pas contenue dans la Proposition~\ref{prop:I-etat-lent-aligne}. Pour une lecture binaire munie de la distance discrète, une telle
borne ne peut pas être uniforme au voisinage de
$\Sigma_{\mathcal P,\rho}$ ; il faut exclure une bande de frontière,
comme dans Proposition~\ref{prop:H-extinction-approchee}.

\paragraph{Changement de résolution.}
Une coordonnée inactive à une résolution grossière peut devenir active
à une résolution plus fine. La suffisance et l'activité sont toujours
relatives à $(\mathfrak B,\rho)$.

\section{Dépendances logiques et usages}
\label{annexe-I-dependances}

La chaîne de dépendance est :
\[
\begin{aligned}
&\text{A-H1 + O1 + invariance alignée + inertie des variables omises}\\
&\qquad\Longrightarrow \text{Proposition~\ref{prop:I-etat-lent-aligne}}\\
&\qquad\Longrightarrow \text{carte }\widehat\Omega_{\mathcal P,\rho}\\
&\qquad\Longrightarrow \text{frontière }\Sigma_{\mathcal P,\rho}.
\end{aligned}
\]

La proposition est utilisée de quatre manières :
\begin{enumerate}
\item elle identifie l'espace minimal actuellement démontré pour
  cartographier la réponse des histoires non complètement éteintes ;
\item elle replace le Théorème~\ref{thm:H-contenu-survivant} comme une section de la géométrie
  lente complète ;
\item elle permet d'interpréter facilitation et anti-facilitation comme
  des déplacements de l'état déposé relativement à
  $\Sigma_{\mathcal P,\rho}$ ;
\item elle prépare la limite de la Proposition~\ref{prop:J-pli-couture} : une modification du mécanisme
  d'exécution peut déformer la frontière même lorsque le dépôt historique
  est maintenu fixe.
\end{enumerate}

\section{Table des notations}
\label{annexe-I-notations}

\begin{center}
\begin{tabular}{p{0.25\textwidth}p{0.62\textwidth}}
\toprule
Notation & Signification \\
\midrule
$z=(p,\chi,\theta,s)$ & état lent aligné \\
$Z$ & domaine déclaré des états lents atteignables \\
$\mathcal M_{\mathrm{al}}$ & sous-espace invariant des états alignés \\
$\iota(z;ssc)$ & injection de l'état lent dans l'état complet regrainé \\
$\widehat\Omega_{\mathcal P,\rho}$ & carte de réponse sur l'état lent aligné \\
$\Omega_{\mathcal P,\rho}$ & restriction à la section éteinte, carte nue de l'Appendice~H \\
$A_{\mathcal P,\rho}$ & ensemble d'appropriation dans $Z$ \\
$\Sigma_{\mathcal P,\rho}$ & frontière relative de l'ensemble d'appropriation \\
$H$ & carte de dépôt d'une famille historique dans $Z$ \\
\bottomrule
\end{tabular}
\end{center}

\section{Conclusion}
\label{annexe-I-conclusion}

La suffisance de l'état lent aligné repose sur une fermeture précise :
les exécutions restent dans le sous-espace aligné, les variables omises
sont inertes et la regraine fixe la composante mécanique. Sous ces
conditions, toute histoire n'agit sur la sonde que par le quadruplet
qu'elle dépose,
\[
(p,\chi,\theta,s).
\]
La carte de réponse nue du régime éteint en est la section
$p=\chi=\theta=0$. La frontière d'appropriation devient ainsi un objet
dans l'état lent complet, et les histoires peuvent être étudiées comme
des trajectoires ou des cartes de dépôt qui la rencontrent.

Le résultat est volontairement local à la famille et à la loi déclarées.
Cette limitation n'est pas une faiblesse : elle indique exactement ce
qui doit être agrandi lorsque l'alignement, l'inertie de $ssc$ ou la
regraine canonique cessent d'être valides.
